\documentclass[10pt]{article}

\usepackage[letterpaper,margin=0.78in]{geometry}
\usepackage{booktabs}
\usepackage{microtype}
\usepackage{tabularx}
\usepackage[table]{xcolor}
\usepackage{hyperref}
\usepackage{amsmath}
\usepackage{xspace}

\definecolor{scopeblue}{HTML}{16324F}
\definecolor{scopeaccent}{HTML}{2D6A8A}
\definecolor{scopewash}{HTML}{EDF4F7}
\definecolor{scopegray}{HTML}{59636E}

\hypersetup{
  colorlinks=true,
  linkcolor=scopeaccent,
  urlcolor=scopeaccent,
  pdftitle={Token Cost and Early-Termination Break-Even Analysis for Scopey},
  pdfauthor={ArchAstro}
}

\newcommand{\scopey}{\textsc{Scopey}\xspace}
\newcommand{\tokens}{\text{tokens}}
\newcommand{\result}[1]{\textcolor{scopeblue}{\textbf{#1}}}
\renewcommand{\arraystretch}{1.14}

\title{\vspace{-1.5em}\textbf{Token Cost and Early-Termination Break-Even Analysis for \scopey}\\
\large A controlled component evaluation of local Qwen3.5 9B Q4}
\author{ArchAstro Engineering}
\date{3 August 2026}

\begin{document}
\maketitle

\begin{abstract}
We extend the \scopey scope-transition benchmark with explicit token accounting.
The experiment compares a no-\scopey control with the selected local summarizer,
Qwen3.5 9B Q4 served by llama.cpp, over 12 synthetic scenarios, 22 scored turns,
and three repetitions. Per repetition, the measured component workload contains
431 main-session user-context tokens, 7,958 \scopey analyzer-input tokens, and
1,613 \scopey-generated tokens. Total analyzer overhead is 9,571 tokens, or 435
tokens per scored turn. Because the component corpus contains no main-agent
generations or tool transcripts, early-termination savings are not observed.
Instead, we report the falsifiable break-even condition: an intervention must
prevent an average of \result{798 future main-session tokens per scenario} to
recover analyzer overhead. If one intervention prevents a 2,500-token suffix in
each scenario, the projected reduction is \result{20,429 tokens per corpus
repetition (67.1\%)}. This projection is a sensitivity result, not a causal
end-to-end measurement.
\end{abstract}

\section{Question and experimental boundary}

The practical question is whether the tokens spent maintaining an explicit
task scope are repaid when \scopey stops an agent from continuing down an
obsolete or off-track trajectory. The accounting separates three quantities:

\begin{enumerate}
  \item \textbf{Main-session tokens} ($M$): cumulative user-request context in
  the component fixture. Main-agent output, reasoning, tool calls, and tool
  results are absent and therefore are not counted.
  \item \textbf{\scopey analyzer input} ($S_i$): tokens submitted to the local
  scope-state transformer, including its instructions, prior scope, recent
  requests, and latest request.
  \item \textbf{\scopey-generated tokens} ($S_o$): completion tokens produced
  by the local model. Total analyzer overhead is $H=S_i+S_o$.
\end{enumerate}

This boundary matters. The evaluation measures summarization cost exactly for
the chosen local runtime, but it does not measure a full coding-agent session.
Consequently, avoided future work is expressed as a parameter rather than
silently inferred from synthetic scope text.

\section{Method}

\subsection{Configurations}

The control exposes only the latest user prompt and performs no model-backed
scope analysis. The treatment uses the selected local Qwen3.5 9B Q4 GGUF model
through a warm llama.cpp server. Decoding uses temperature zero, seed 42, a
256-token completion limit, and validated JSON normalized into \scopey's scope
format. The runtime revision is
\texttt{fe2adf0e722f30f5\allowbreak{}295fdec8a0f1dc78\allowbreak{}8f7498bc}; the
model SHA-256 is
\texttt{d784ce9eda1a5a7b\allowbreak{}51e8f705a9e63108\allowbreak{}44bf4f173654d115\allowbreak{}823c775fdea56d43}.

The 12 cases exercise additions, narrowing, subtraction, explicit and implicit
replacement, read-only queries, administration, combined mutations,
constraint preservation, non-invention, and reference resolution. Each run
contains 22 scored scope transitions. Three repetitions produce 66 samples per
configuration and 132 samples overall.

\subsection{Counting and projection}

Qwen tokenization is used for main-session fixture text. The llama.cpp
OpenAI-compatible response supplies analyzer prompt and completion usage, so
the \scopey counts are runtime-reported rather than byte estimates. Counts in
Tables~\ref{tab:aggregate}--\ref{tab:scenarios} are normalized to one corpus
repetition.

Let $N=12$ scenarios and let $A$ be the future main-session tokens prevented by
one successful early termination in each scenario. The modeled totals are

\begin{align}
T_{\mathrm{no\ scopey}}(A) &= M + NA,\\
T_{\mathrm{scopey}} &= M + H,\\
\Delta(A) &= NA-H.
\end{align}

Positive $\Delta$ indicates net token savings. The corpus breaks even at
$A=H/N=9{,}571/12=797.6$ tokens per scenario. Scenario-specific break-even
values use each scenario's measured analyzer overhead.

\section{Results}

\begin{table}[ht]
\centering
\caption{Aggregate token accounting per 12-scenario corpus repetition.}
\label{tab:aggregate}
\begin{tabular}{lrr}
\toprule
Token bucket & No \scopey & Local \scopey \\
\midrule
Main-session user context ($M$) & 431 & 431 \\
Analyzer input ($S_i$) & 0 & 7,958 \\
\scopey-generated output ($S_o$) & 0 & 1,613 \\
Total analyzer overhead ($H$) & 0 & \result{9,571} \\
Observed component total & 431 & 10,002 \\
\bottomrule
\end{tabular}
\end{table}

Analyzer input accounts for 83.1\% of \scopey overhead and generated output for
16.9\%. The model generates 73.3 tokens per scored turn on average; full
analyzer cost is 435.0 tokens per turn. The no-\scopey control has no analyzer
cost. Without an avoided future trajectory, it is necessarily cheaper in token
terms; \scopey's value appears only when a correction prevents enough later
main-session work.

\clearpage
\begin{table}[ht]
\centering
\caption{Per-scenario token cost and break-even, averaged over three repetitions.
``Main'' is cumulative user-request context. ``Generated'' is model output.}
\label{tab:scenarios}
\small
\begin{tabularx}{\textwidth}{@{}Xrrrrrr@{}}
\toprule
Scenario & Turns & Main & Input & Gen. & Cost & Net at 2.5k \\
\midrule
Initial concrete addition & 1 & 15 & 352 & 65 & 417 & 2,083 \\
Follow-up preserving work & 2 & 32 & 707 & 120 & 827 & 1,673 \\
Narrow target & 2 & 42 & 738 & 147 & 885 & 1,615 \\
Subtract requirement & 2 & 44 & 738 & 143 & 881 & 1,619 \\
Explicit task replacement & 2 & 47 & 726 & 157 & 883 & 1,617 \\
Implicit topic switch & 2 & 40 & 717 & 140 & 857 & 1,643 \\
Query preserves work & 2 & 38 & 715 & 147 & 862 & 1,638 \\
Administration preserves work & 2 & 31 & 713 & 141 & 854 & 1,646 \\
Combined subtract/add & 2 & 47 & 737 & 158 & 895 & 1,605 \\
Cross-platform constraints & 2 & 51 & 757 & 201 & 958 & 1,542 \\
No invented administration & 1 & 10 & 347 & 53 & 400 & 2,100 \\
Reference resolution & 2 & 34 & 711 & 141 & 852 & 1,648 \\
\midrule
\textbf{Corpus total} & \textbf{22} & \textbf{431} & \textbf{7,958} &
\textbf{1,613} & \textbf{9,571} & \textbf{20,429} \\
\bottomrule
\end{tabularx}
\end{table}

\subsection{Early-termination sensitivity}

Table~\ref{tab:sensitivity} makes the counterfactual assumption visible. At
zero avoided tokens, \scopey adds its full 9,571-token analyzer cost. Savings
turn positive near 798 tokens per scenario. At the requested illustrative
2,500-token suffix, the no-\scopey continuation would consume 30,431 tokens,
whereas the terminated \scopey path consumes 10,002, a 67.1\% reduction.

\begin{table}[ht]
\centering
\caption{Sensitivity to assumed future main-session tokens prevented per scenario.}
\label{tab:sensitivity}
\begin{tabular}{rrr}
\toprule
Avoided per scenario ($A$) & Net corpus savings $\Delta(A)$ & Reduction \\
\midrule
0 & -9,571 & -2,220.6\% \\
250 & -6,571 & -191.5\% \\
500 & -3,571 & -55.5\% \\
800 & 29 & 0.3\% \\
1,000 & 2,429 & 19.5\% \\
2,500 & \result{20,429} & \result{67.1\%} \\
5,000 & 50,429 & 83.4\% \\
\bottomrule
\end{tabular}
\end{table}

The local model retained the earlier quality result during this run: 100\%
format compliance, 100\% required-concept recall, 100\% forbidden-concept
rejection, zero call failures, and 59.1\% exact transition-label agreement.
Median latency was 1,190 ms. The lower transition-label score reflects a
coarser operation taxonomy; the active scope body consumed by the product
remained correct on all measured concepts.

\section{Threats to validity}

\begin{itemize}
  \item \textbf{No causal termination measurement.} The 2,500-token suffix is a
  sensitivity parameter. No main agent was actually stopped in this component
  experiment.
  \item \textbf{Partial main-session accounting.} The 431-token main bucket
  contains user-request context only. It omits system prompts, assistant output,
  reasoning, cache reads/writes, and tool payloads.
  \item \textbf{One tokenizer.} Counts use the selected Qwen tokenizer. Cloud
  providers may tokenize identical text differently.
  \item \textbf{Synthetic cases.} The corpus is deterministic and auditable but
  is smaller and cleaner than long production coding sessions.
  \item \textbf{Intervention frequency.} The projection assumes one successful
  termination per scenario. Real prevalence and false-positive rates determine
  realized savings.
\end{itemize}

\section{Conclusion and next experiment}

Local scope summarization is not token-free: the selected configuration costs
about 435 analyzer tokens per scored turn and 798 tokens per scenario. It is
token-positive only when a correction prevents more future main-session work
than that overhead. Under a 2,500-token prevented-suffix scenario, the modeled
benefit is substantial, but it remains projected.

The next evaluation should run paired disposable coding tasks with matched
model, prompt, seed, repository state, and stopping rule. It must capture
provider-reported input, output, cached, and reasoning tokens at every main-agent
step; record the precise intervention token index; continue the no-\scopey arm
to completion; and score task success plus false interventions. That experiment
would replace $A$ with observed counterfactual trajectory differences and turn
this break-even analysis into an end-to-end causal estimate.

\section*{Reproducibility}

The evaluator, case corpus, run command, and machine-readable schema live under
\path{eval/}. The source paper is
\path{eval/paper/scopey-token-accounting.tex}. The selected model and runtime
are documented in \path{eval/LOCAL_MODEL_DECISION.md}. Runtime and model
projects: \url{https://github.com/ggml-org/llama.cpp} and
\url{https://huggingface.co/Qwen/Qwen3.5-9B}.

\end{document}
